\chapter*{Введение: зачем этот фундамент}
\addcontentsline{toc}{chapter}{Введение: зачем этот фундамент}
\markboth{Введение}{Введение}

Водород, попавший в оболочку твэла из циркониевого сплава, при охлаждении выпадает в виде тонких пластинок
гидрида. Обычно они лежат вдоль окружности трубы и почти не мешают ей держать давление. Но если труба остывает под
растягивающим окружным напряжением, часть пластинок встаёт поперёк --- по радиусу (рис.~\ref{fig:f0_tube}).
Радиальные гидриды режут стенку поперёк и делают металл хрупким. Модель, ради которой написан этот текст,
предсказывает, сколько гидридов переориентируется при заданных водороде, температуре, напряжении и структуре
металла.

\fig{f0_tube}{Переориентация гидридов в сечении трубы (схема). Слева --- до опыта: пластинки гидрида (синие)
лежат вдоль окружности. Справа --- после охлаждения под окружным напряжением $\sigma_\theta$: часть пластинок
встала по радиусу (красные).}

Чтобы такую модель построить или хотя бы понять, нужно одновременно держать в голове несколько областей:
термодинамику растворов, диффузию, кристаллографию, теорию упругости, теорию зарождения и пластичность.
Каждая из них --- отдельный учебник. Здесь собрано только то, что реально работает в задаче о гидридах,
и объяснено так, чтобы было видно, \emph{почему} формулы такие, а не только \emph{какие} они.

\section*{Как устроен текст}

\begin{center}
\small
\begin{tabularx}{\linewidth}{>{\raggedright\arraybackslash}p{4.6cm}X>{\centering\arraybackslash}p{1.4cm}}
\toprule
что есть в модели & что нужно понимать & глава \\
\midrule
пределы растворимости водорода, гистерезис растворения и выпадения & энергия Гиббса, химический потенциал, растворы, общая касательная & 1, 2 \\
перенос водорода между гидридами, к растянутым местам & законы Фика, химический потенциал как движущая сила, эффект Горского & 3 \\
форма и плоскость залегания пластинок, межфазные границы & решётка ГПУ, гидриды циркония, когерентность, дислокации & 4 \\
энергия пластинки, поле напряжений вокруг неё, работа нагрузки & собственная деформация, задача Эшелби, метод Фурье & 5 \\
когда и где появляется пластинка, как растёт & классическая теория зарождения, неоднородное зарождение, рост, кинетика & 6 \\
пластическая аккомодация, память, текстура трубы & текучесть, дислокационная пластичность, текстура и факторы Кернса & 7 \\
\bottomrule
\end{tabularx}
\end{center}

Главы можно читать по порядку: каждая опирается на предыдущие. В каждой есть блоки трёх видов.
\begin{itemize}
\item \textbf{\color{circ}Интуиция} --- что происходит «на пальцах», до формул.
\item \textbf{\color{load}Главная формула} --- то, что стоит запомнить или уметь вывести.
\item \textbf{\color{hydr}Для гидридов} --- как общий результат применяется к водороду в цирконии, с числами.
\end{itemize}
В конце глав --- упражнения; ответы и подсказки собраны в приложении Б. Справочные величины для системы Zr--H
с источниками --- в приложении А. Все рисунки посчитаны по формулам этого текста (скрипты лежат в
\texttt{figs\_src/}); числа на них можно проверить самостоятельно.

\section*{Единицы и масштабы}

В задаче о гидридах встречаются энергии «на атом», «на моль» и «на кубический метр»; их полезно уметь переводить
друг в друга с ходу.
\begin{itemize}
\item Тепловая энергия: $k_BT$ на атом или $RT$ на моль. При 300\degC{} ($T = 573$~К)
  $RT \approx 4.8$~кДж/моль, $k_BT \approx 0.049$~эВ. Это мера «тепловых толчков»: барьер в $20RT$ при этой
  температуре преодолевается редко, в $2RT$ --- постоянно.
\item Концентрацию водорода в цирконии обычно дают в \emph{ppm по массе} (wppm, млн$^{-1}$): микрограммы водорода
  на грамм металла. Атомная доля $x_{\mathrm H} \approx C_{\mathrm{ppm}} \cdot 91.22 / (1.008\cdot 10^6)$, то есть
  100~ppm $\approx 0.9$~ат.\,\%.
\item Энергия на единицу объёма: 1~МДж/м$^3$ = 1~МПа. Напряжение --- это и есть плотность энергии. Поэтому
  работа нагрузки 100~МПа на деформации 7\,\% равна 7~МДж/м$^3$, и её можно прямо сравнивать с упругой энергией
  гидрида в тех же единицах (глава~5).
\item Объём гидрида на моль водорода около $10^{-5}$~м$^3$. Поэтому 1~МДж/м$^3$ гидрида $\approx 10$~Дж на моль водорода.
\end{itemize}
